\documentclass{article}
\usepackage[margin=1.5cm]{geometry}
\usepackage{amsmath}
\usepackage{enumitem}

\setlist{nosep,leftmargin=*}

\begin{document}
\twocolumn

\title{Chapter 7 Warm-Up: Importing Functions from Modules}
\author{Prof. Jordan C. Hanson}
\date{}

\maketitle

\section{Creating a Module}

\begin{enumerate}

\item A Python \textbf{module} is a \texttt{.py} file containing
code that can be imported into another program.

Create a new file called

\begin{verbatim}
math_tools.py
\end{verbatim}

and place the following functions inside it:

\begin{verbatim}
def square(x):
    return x * x

def cube(x):
    return x * x * x

def average(x, y):
    return (x + y) / 2
\end{verbatim}

Save the file.  The functions are now organized inside the
module \texttt{math\_tools}.

\end{enumerate}


\section{Uploading to Google Colab}

\begin{enumerate}

\item Open a new Google Colab notebook.

Upload \texttt{math\_tools.py} to the Colab runtime:

\begin{enumerate}[label=\arabic*.]
\item Click the \textbf{Files} folder icon on the left side.
\item Click the \textbf{Upload} icon in the Files panel.
\item Select \texttt{math\_tools.py} from your computer.
\item Confirm that \texttt{math\_tools.py} appears in the
      Colab file list.
\end{enumerate}

The module must be in the current Colab runtime before Python can
import it.

\end{enumerate}


\section{Importing the Module}

\begin{enumerate}

\item In a Colab code cell, import the complete module:

\begin{verbatim}
import math_tools
\end{verbatim}

Functions inside the module are accessed using the module name
followed by a period:

\begin{verbatim}
print(math_tools.square(5))
print(math_tools.cube(3))
print(math_tools.average(8, 12))
\end{verbatim}

Before running the cell, predict the three outputs.

\begin{itemize}

\item What does the name \texttt{math\_tools} refer to?

\item Which part of

\begin{verbatim}
math_tools.square(5)
\end{verbatim}

is the module name?  Which part is the function name?

\end{itemize}

\end{enumerate}


\section{Importing a Function}

\begin{enumerate}

\item Python can also import an individual function:

\begin{verbatim}
from math_tools import square
\end{verbatim}

Now call it without the module name:

\begin{verbatim}
print(square(7))
\end{verbatim}

Compare this with

\begin{verbatim}
import math_tools
print(math_tools.square(7))
\end{verbatim}

Both call the same function, but they use different import
statements.

\end{enumerate}


\section{Code-Lab Exercise}

\begin{enumerate}

\item Add the following function to
\texttt{math\_tools.py}:

\begin{verbatim}
def classify_score(score):
    if score >= 4:
        return "Calculus"
    elif score >= 2:
        return "Trigonometry"
    else:
        return "Algebra"
\end{verbatim}

Upload the updated file to Colab.  Then import the function:

\begin{verbatim}
from math_tools import classify_score
\end{verbatim}

Test it using

\begin{verbatim}
scores = [0, 2, 3, 5]

for score in scores:
    course = classify_score(score)
    print(score, course)
\end{verbatim}

\begin{itemize}

\item Predict the four course assignments before running the code.

\item Which file \textbf{defines} the function?

\item Which file or notebook \textbf{calls} the function?

\item Add one more score to the list and verify that the imported
function still performs the placement correctly.

\end{itemize}

\end{enumerate}


\section{Important}

Files uploaded directly to a Colab runtime are temporary.  If the
runtime is restarted or disconnected, you may need to upload your
module again.

\end{document}